% Ad Atticum 4.14 --- headnote
% ad-atticum-04-14, mid-May 54 BC, Cumae

\textit{Cicero to Atticus, written from his Cumanum just after 10
May 54 BC. Vestorius --- the Puteoli banker who is the brothers'
shared correspondent on the bay of Naples --- has reported
Atticus's late departure from Rome and the illness behind it.
Cicero's main request is for the run of Atticus's library while
the owner is away: ``especially Varro's,'' he says. The work
in his hands is almost certainly \textit{De Re Publica}, drafted
through 54 BC and published in 51, for which Varro's antiquarian
materials would be exactly what he needed.}

\textit{Section 2 is the standard Cicero-on-Atticus's-letters
formula: a request for any news, especially from Quintus (now in
Caesar's camp in Gaul) and from Caesar himself; on the elections;
on the commonwealth in any form. ``You tend to scent these things
wittily.'' Even an empty letter is welcome. Closing with a greeting
to Dionysius and a wish for Atticus's safe return.}
